% !TEX root = ../main_zh.tex
\chapter{索洛增长模型}
\label{ch:solow}

为什么今天一个普通美国人的富裕程度，大约是两百年前普通人的五十倍？为什么今天有些国家的富裕程度是另一些国家的二十倍？这些关于\textbf{长期}的问题——关于生活水平在数十年乃至数代人之间的缓慢漂移——正是增长理论的研究领域，而它们的意义极为重大：正是人均产出的持续增长，提升了消费、延长了预期寿命，并通过把经济总量这块蛋糕做大，使得即便在不平等程度很高的地方，再分配与社会稳定也成为可能。我们后面要研究的经济周期，也就是短期波动，最好被理解为围绕这条上升趋势线的偏离；用一句有用的口号来说，周期=增长+冲击。要理解周期，我们首先需要一套关于趋势的理论。

本章发展第一个、也是最简单的这样一套理论：\textbf{索洛增长模型}（Solow growth model）。在引入它之前，我们先停下来搭好所有宏观动态模型共享的一般框架：一条时间线、一个概括经济当前所处状态的状态变量、经济所选择的控制变量，以及一条把状态推向未来的运动方程。索洛模型是这个框架最精简的一个实例。它标志性的简化处理——也是把它与第~\ref{ch:neoclassical}~章新古典模型区分开来的那唯一一条假设——在于储蓄率是外生固定的，而非由优化的家庭选择得到的。我们将搭起基准模型、求出它的稳态、追问哪个稳态是最优的（黄金律）、刻画政策改变储蓄率后经济如何调整，最后把模型扩展到容纳人口增长与技术进步。这个模型不能解释一切，但它能解释相当多的东西，并且它把增长理论其余部分所要回答的问题组织了起来。

\section{宏观经济学的动态框架}

微观经济学在很大程度上是静态的：消费者选择一个消费束，厂商选择一个产量，我们再比较均衡。而宏观增长在本质上是\textbf{动态}的。其核心对象不是单个选择，而是一串通过时间相连的选择，把它们连起来的装置就是\textbf{资本}的积累。因此，每一个动态宏观模型都从铺设一条时间线开始。

\state{解决宏观问题一定要先有一条时间线}{
要分析任何宏观增长问题，先写下一条以时期 $t = 1, 2, \dots$ 为标号的时间线。每一期，经济体继承一笔资本存量，做出生产出产出、并决定带入下一期多少资本的种种选择，如此循环往复。
}

资本存量扮演着特殊的角色。在第 $t$ 期之初，存量 $K_t$ 是固定的——它就是过去的投资所遗留下来的那个数——所以从第 $t$ 期的视角看，它是\textbf{给定}的，是今天的经济体无法改变的一个数。这种变量被称为\textbf{状态变量}（state variable）：它概括了经济体为了决定眼下该怎么做、所需要知道的关于自身过去的一切。面对给定的状态，经济体选择如何把当期产出在消费与储蓄之间分配。这些选择就是\textbf{控制变量}（control variable）：消费 $C_t$ 与储蓄 $S_t$ 可以自由设定，而一旦设定，下一期的资本存量就被确定了。决定下一期存量的那条规则，就是\textbf{资本的运动方程}（law of motion of capital），它刻画了存量如何随时间积累。

\defn{状态变量与控制变量}{
在一个动态模型中，\textbf{状态变量}是一个在每期之初已被预先决定、并概括了经济体相关历史的量——这里就是资本存量 $K_t$。\textbf{控制变量}是经济体在当期内自由选择的量——这里就是消费 $C_t$ 与储蓄 $S_t$。状态约束着控制，控制又决定下一期的状态。
}

为把这些想法落到实处，考虑一个去掉了家庭与厂商、只展示支配任何经济体的会计恒等式的两期例子。产出由资本生产，$Y_i = F(K_i)$；第一期产出在消费与储蓄之间分配，$C_1 + S_1 = Y_1$；储蓄被投资出去，$I_1 = S_1$；资本以 $\delta$ 的速率折旧、又由投资补充，故 $K_2 = (1-\delta)K_1 + I_1$；第二期家庭消费其产出再加上储蓄的总回报，$C_2 = Y_2 + (1+r)S_1$。把这些汇总起来，
\[
\ba
C_1 + S_1 &= Y_1, \\
C_2 &= Y_2 + (1+r) S_1, \\
K_2 &= (1-\delta)K_1 + I_1, \qquad I_1 = S_1, \\
Y_i &= F(K_i).
\ea
\]
这些就是经济体机械运行所遵循的规律；它们不涉及任何优化。这恰恰是索洛模型工作的层次——它给产出如何分配施加一条固定规则，然后任凭会计恒等式把经济体推向未来。而第~\ref{ch:neoclassical}~章的新古典模型则会加上微观基础，从一个最大化终身效用的家庭、以及在一般均衡中出清的厂商与市场出发，推导出消费-储蓄的分配。具体而言，优化的家庭求解
\[
\max \ \sum_{t=1}^{\infty} \beta^{t-1} U(C_t),
\]
约束条件就是同样那几条运动方程。但无论我们是否加上优化，结构都是一样的：每一期 $K_t$ 是状态，$C_t$ 与 $S_t$ 是控制。

这一结构有一个深刻的后果。站在第 $t$ 期、手握状态 $K_t$，每选定一个 $C_t$ 就钉死了一个 $S_t$，进而钉死了 $K_{t+1}$。于是从今天的状态到明天的状态的映射，可以写成单独一个函数，
\[
K_{t+1} = g(K_t),
\]
称为\textbf{政策函数}（policy function）。一期接一期地顺着政策函数走，就描画出整个经济体的最优路径。这个映射有两个值得点名的特征。第一，下一期的状态只依赖于这一期的状态，而不依赖经济体是经由怎样的整段历史走到这里的；这个问题具有\textbf{马尔可夫}（Markov）性质。第二，每个日期所面对的决策问题结构完全相同——不同的只是所继承的状态取值——所以这个问题是\textbf{递归的}（recursive）。这条乍看令人生畏的无穷期选择序列，于是坍缩成同一个问题、以不同起点被反复求解。正是这种递归性使得动态宏观可处理，它支撑起本书中的每一个增长模型。

\section{索洛模型：基本假设}

现在我们把这个框架专门化到索洛的经济中。这个模型作出六条贯穿始终的假设，我们把它们集中放在一处。

\asm{索洛模型}{
\begin{enumerate}
    \item \textbf{总量生产函数。}产出由资本与劳动生产，$Y = F(K, L)$。劳动力 $L$ 保持不变；唯一会变动的量是资本 $K$。
    \item \textbf{规模报酬不变（CRS）。}$F$ 是一次齐次的：把两种投入都乘以任一倍数 $\lambda$，产出也乘以 $\lambda$。
    \item \textbf{边际报酬递减与稻田条件。}每种投入都是有生产力的，但报酬递减，
    \[
    F_K > 0, \quad F_L > 0, \qquad F_{KK} < 0, \quad F_{LL} < 0,
    \]
    且当某种投入趋于消失时其边际产出变得无界，
    \[
    \lim_{K \to 0} F_K(K, L) = \infty, \qquad \lim_{L \to 0} F_L(K, L) = \infty.
    \]
    \item \textbf{无政府的封闭经济。}净出口与政府支出均为零，$\mathrm{NX} = G = 0$。
    \item \textbf{储蓄率外生且不变。}产出中一个固定份额 $s \in (0,1)$ 被储蓄起来，$I = sY$，其中 $s$ 由模型之外给定。
    \item \textbf{资本积累。}资本以 $\delta$ 的速率折旧、又由投资补充，
    \[
    K' = (1-\delta) K + I,
    \]
    其中撇号表示下一期。
\end{enumerate}
}

在动用它们之前，其中有两条假设值得评说。

规模报酬不变（假设 2）正是让我们能够在总量经济与一个代表性的“劳均”描述之间自由穿梭的东西。在 CRS 下，一家厂商被拆分或合并都不改变其生产率，所以我们既可以把整个社会看作一家雇用 $L$ 个工人与 $K$ 台机器的巨型厂商，也可以等价地看作单独一个代表性工人操作着 $k := K/L$ 单位资本。在齐次性性质中取 $\lambda = 1/L$，
\[
\frac{Y}{L} = F\!\pr{\frac{K}{L}, 1} =: f(k),
\]
这就定义了\textbf{集约型}（劳均）生产函数 $y = f(k)$。稻田条件（假设 3）保证了内点解：因为当资本稀缺时其边际产出极大、当资本充裕时其边际产出极小，经济体总是被从 $k = 0$ 与 $k = \infty$ 这两个角点拉开，拉向一个合理的内点稳态。

封闭经济假设（假设 4）让我们能把储蓄与投资划上等号。在 $\mathrm{NX} = G = 0$ 下，产出恒等式读作 $Y = C + I$，而收入被分为消费与储蓄，$Y = C + S$；两式相减便得到封闭经济的基本条件
\[
S = I.
\]
化为劳均形式即 $y = c + i$，其中 $i = sy$、$c = (1-s)y$。

\state{索洛之所以为“索洛”}{
这个模型的简单之处——也是把它与新古典模型区分开来的那一条假设——在于\textbf{储蓄率 $s$ 是外生且不变的}。家庭并不通过权衡当前效用与未来效用来选择储蓄多少；收入中一个固定份额被机械地储蓄起来。第~\ref{ch:neoclassical}~章的新古典模型则用一个最优的、最大化效用的消费-储蓄选择来取代这一点。正是在这个确切的意义上，索洛描述的是一个储蓄规则被硬性写死的计划经济。
}

\section{基准模型：无人口增长}

我们从最简单的情形入手，即劳动力恒定。由于 $L$ 不变，劳均的运动方程与总量形式具有相同的样子。把 $K' = (1-\delta)K + I$ 除以常数 $L$，
\[
k' = (1-\delta) k + i.
\]
于是劳均资本存量从一期到下一期的变化为
\begin{equation}
\Delta k = k' - k = i - \delta k = s\, y - \delta k = s\, f(k) - \delta k.
\label{eq:solow-lom}
\end{equation}
式~\eqref{eq:solow-lom} 是索洛模型的心脏。它是劳均资本的运动方程，读懂它就是读懂整个模型。当劳均总投资 $s f(k)$ 超过折旧损耗的资本 $\delta k$ 时，劳均资本上升；不足时则下降。其动态完全是跨期的：今天在消费与投资之间的分配决定了明天的资本存量，进而决定明天的产出，如此往复。这正是增长与微观经济学静态比较的对照之处——在宏观增长里，戏份全在运动方程上。

\subsection{稳态}

当劳均资本不再变化时，经济体便安顿下来。在式~\eqref{eq:solow-lom} 中令 $\Delta k = 0$，得到\textbf{稳态}条件
\begin{equation}
s\, f(k^*) = \delta k^*.
\label{eq:solow-ss}
\end{equation}

\defn{稳态}{
索洛经济的一个\textbf{稳态}（steady state）是劳均资本的一个水平 $k^*$，在该水平上总投资恰好补足折旧掉的资本，$s f(k^*) = \delta k^*$，从而 $\Delta k = 0$。在稳态处，每一个劳均变量——$k^*$、$y^* = f(k^*)$、$c^* = (1-s) f(k^*)$、$i^* = s f(k^*)$——都随时间保持不变。
}

稳态用图来看最为直观。图~\ref{fig:macro-2} 对着劳均资本 $k$ 画出三条曲线：生产函数 $f(k)$，因报酬递减而递增却上凸；投资曲线 $s f(k)$，与前者同形，只是按 $s$ 缩小了；以及折旧线 $\delta k$，一条过原点的射线。因为在原点附近 $s f(k)$ 起步时高于 $\delta k$（稻田条件使得投资在那里斜率无穷），而最终又被这条直线反超（报酬递减把 $f$ 压平），二者恰好在 $k^*$ 处相交一次。

\begin{figure}[h]
\centering
\begin{tikzpicture}[scale=1.0]
  % f(k) = A*k^a ; sf(k) = s*A*k^a ; delta-line = d*k   (a<1/2: strongly concave)
  \def\A{2.60}      % production scale
  \def\s{0.55}      % savings rate
  \def\a{0.33}      % Cobb-Douglas exponent (small => sharply concave f(k))
  \def\d{0.95}      % depreciation slope (steep enough that delta*k crosses f(k) in-frame)
  \def\xmax{8.0}    % rightmost k drawn
  % steady state: s*A*(k*)^a = d*k*  =>  k* = (s*A/d)^(1/(1-a))
  \pgfmathsetmacro{\ks}{(\s*\A/\d)^(1/(1-\a))}   % k*
  \pgfmathsetmacro{\ys}{\d*\ks}                  % height at k* (on both line and sf)
  % point where delta*k overtakes f(k):  k = (A/d)^(1/(1-a))
  \pgfmathsetmacro{\kc}{(\A/\d)^(1/(1-\a))}      % crossing of delta*k and f(k)
  % axes
  \draw[-Latex] (0,0) -- (9.0,0) node[right] {$k$};
  \draw[-Latex] (0,0) -- (0,6.2) node[above] {$y$};
  % production function y=f(k)  (top concave curve)
  \draw[line width=1.1pt,smooth,samples=140,domain=0:\xmax]
        plot (\x,{\A*(\x)^\a});
  \node[above] at (2.6,{\A*(2.6)^\a+0.18}) {$y=f(k)$};
  % savings / investment curve sy=sf(k)  (lower concave copy)
  \draw[line width=1.1pt,smooth,samples=140,domain=0:\xmax]
        plot (\x,{\s*\A*(\x)^\a});
  \node[right] at (\xmax,{\s*\A*(\xmax)^\a}) {$sy=sf(k)$};
  % depreciation line delta*k  (straight through origin; visibly crosses f(k) past \kc)
  \pgfmathsetmacro{\xline}{\kc+1.2}              % continue the line clearly above f(k)
  \draw[line width=1.1pt] (0,0) -- (\xline,{\d*\xline});
  \node[above right=-1pt and 0pt] at (\xline,{\d*\xline}) {$\delta k$};
  % steady state marked with a star
  \node[star,star points=5,star point ratio=2.2,fill=blue!70,draw=blue!70,
        inner sep=0pt,minimum size=11pt] at (\ks,\ys) {};
  % guide line to k* on the axis
  \draw[densely dashed] (\ks,\ys) -- (\ks,0) node[below] {$k^{*}$};
\end{tikzpicture}

\caption{基准索洛图。经济体收敛到稳态 $k^*$，在那里投资曲线 $s f(k)$ 与折旧线 $\delta k$ 相交；在 $k^*$ 左侧投资超过折旧、$k$ 上升，在右侧则下降。}
\label{fig:macro-2}
\end{figure}

这个稳态是\textbf{稳定的}：无论经济体被推离 $k^*$ 多远，它都会回来。在 $k^*$ 左侧，投资 $s f(k)$ 位于折旧 $\delta k$ 之上，故 $\Delta k > 0$，资本朝 $k^*$ 增长；在右侧，折旧占上风，$\Delta k < 0$，资本又缩回 $k^*$。这一收敛背后的经济学，正是把生产函数压弯的那同一个报酬递减：当资本稀缺时，每多一台机器都极有生产力、轻易就赚回了自身的折旧，于是存量不断累积；随着资本积累，边际产出下降，直到一台新机器的产出恰好只够补上整个存量的折旧，增长便就此停止。

\state{单靠资本积累的增长终将停止}{
在基准索洛模型中，劳均资本的积累最终会停下来。由于资本的边际产出递减、而折旧却线性增长，存在一个有限的 $k^*$ 使二者相抵，经济体便停滞在那里。单靠资本深化，其本身无法维持增长。
}

这一结论带出一个关于收敛的著名推论。两个技术、储蓄率与折旧率相同的经济体，共享同一个 $k^*$，因而无论起点在哪里，都会收敛到相同的劳均产出水平。一个起步时资本很少的后发穷国，最初增长很快——它稀缺的资本赚取高回报——但随着逼近共同的稳态而放慢。早期发展在很大程度上是一个资本积累的故事，但正因为这台发动机会熄火，各国之间任何\textbf{持久}的收入差异都必定来自资本以外的某种东西。于是基准模型过于简单，无法解释我们在世界上观察到的那些持续不散的差距；找出它缺了什么，将驱动下文的各项扩展。

\section{黄金律}

更高的储蓄率会带来更高的稳态资本存量 $k^*$——可是，更富的稳态就是更幸福的稳态吗？这个经济体里的福利是用消费而非资本来衡量的，从这个角度看答案并不显而易见。劳均稳态消费为
\[
c^* = (1-s) f(k^*) = f(k^*) - \delta k^*,
\]
其中第二个等号用了稳态条件 $s f(k^*) = \delta k^*$，把储蓄换成了投资。这样写出来，$c^*$ 就是生产函数 $f(k^*)$ 与折旧线 $\delta k^*$ 之间的竖直缺口：所生产的一切之中，没被折旧吃掉的那部分可供消费。

提高 $s$ 对这个缺口有两个相反的影响。一方面，它把经济体推向更大的 $k^*$，在那里经济体生产得更多、$f(k^*)$ 更高、可分配的东西也更多。另一方面，维持更大的资本存量需要把更多产出仅仅用于补足折旧 $\delta k^*$，这会侵蚀消费。在两个极端之间——在什么都不储蓄（$k^* = 0$，没有产出）与储蓄一切（$c^* = 0$，全部产出都重新投资）之间——的某处，消费达到最大。那个最优的稳态就是\textbf{黄金律}（Golden Rule）。

\thm{黄金律条件}{
在所有稳态中，劳均消费在满足
\[
f'(k_g^*) = \delta
\]
的资本存量 $k_g^*$ 处达到最大。黄金律就是带来这一稳态的那个储蓄率；把这一最优条件与稳态条件 $s f(k_g^*) = \delta k_g^*$ 联立，
\[
s_g = \frac{\delta k_g^*}{f(k_g^*)}.
\]
}

\pf{
稳态消费 $c^*(k^*) = f(k^*) - \delta k^*$ 是储蓄率所挑选的那个稳态资本存量的函数。对 $k^*$ 求最大，
\[
\frac{\d c^*}{\d k^*} = f'(k^*) - \delta = 0
\quad\Longrightarrow\quad
f'(k^*) = \delta.
\]
从几何上看这是一个切线问题：消费是 $f(k^*)$ 与 $\delta k^*$ 之间的缺口，而这个缺口在生产函数斜率等于折旧线斜率 $\delta$ 之处最宽。$f$ 的凹性保证了这是一个最大值。把稳态条件 $s f(k_g^*) = \delta k_g^*$ 在 $k_g^*$ 处读出来，就得到 $s_g$ 的表达式。
}

图~\ref{fig:macro-3} 把这一几何关系画得明明白白。稳态消费是 $f(k^*)$ 与 $\delta k^*$ 之间的竖直距离，而这段距离在 $k_g^*$ 点处最大——在那里生产函数与折旧线平行，也就是 $f'(k_g^*) = \delta$，资本的边际产出等于折旧率。

\begin{figure}[h]
\centering
\begin{tikzpicture}[scale=1.0]
  % --- parameters of the production function f(k) = A*sqrt(k) ---
  % choose A and delta so that the break-even line delta*k crosses f(k)
  % inside the frame, with the golden-rule capital k_g = k^*/4 to its left.
  \def\A{1.451}           % scale of f(k):  crossing at k^*=7.6
  \def\del{0.526}         % slope of the depreciation / break-even line delta*k
  % steady state where f(k)=delta*k:  k^* = (A/delta)^2
  \pgfmathsetmacro{\kstar}{(\A/\del)^2}
  % golden rule: f'(k_g)=A/(2 sqrt(k_g))=delta  =>  k_g = (A/(2 delta))^2 = k^*/4
  \pgfmathsetmacro{\kg}{(\A/(2*\del))^2}
  \pgfmathsetmacro{\fkg}{\A*sqrt(\kg)}     % height of f at k_g
  \def\kmax{8.6}

  % --- axes ---
  \draw[-Latex] (0,0) -- (\kmax,0) node[right] {$k$};
  \draw[-Latex] (0,0) -- (0,5.6) node[above] {output};

  % --- depreciation / break-even line delta*k (black, per house style) ---
  % extend slightly beyond the curve's end so it finishes visibly higher than f(k)
  \draw[line width=1.0pt] (0,0) -- (\kmax,{\del*\kmax}) ;
  \node at ({\kmax-0.85},{\del*(\kmax-0.85)+0.40}) {$\delta k$};

  % --- production function f(k) = A sqrt(k)  (primary, red as in original) ---
  \draw[red,line width=1.5pt,smooth,samples=200,domain=0:\kmax]
        plot (\x,{\A*sqrt(\x)});
  \node[red] at ({\kmax+0.32},{\A*sqrt(\kmax)-0.12}) {$f(k)$};

  % --- tangent line to f(k) at the golden-rule point ---
  % the golden rule is exactly where f'(k_g)=delta, so this tangent has slope
  % EXACTLY delta and is therefore parallel to the delta*k line above.
  % line: y = fkg + del*(x - kg), drawn long & symmetric so the parallelism is obvious.
  % draw the tangent across most of the frame so its slope is visibly identical to delta*k
  \pgfmathsetmacro{\txa}{0.0}
  \pgfmathsetmacro{\txb}{\kmax-1.7}
  \draw[line width=1.0pt] ({\txa},{\fkg+\del*(\txa-\kg)}) -- ({\txb},{\fkg+\del*(\txb-\kg)});
  \node[align=center] at ({\kg+1.15},{\fkg+\del*1.15+0.42}) {$MPK = \delta$};

  % --- golden-rule point on f(k) and dashed guide to the axis ---
  \fill (\kg,\fkg) circle (1.4pt);
  \draw[densely dashed] (\kg,\fkg) -- (\kg,0) node[below] {$k_g^{*}$};
\end{tikzpicture}

\caption{黄金律。稳态消费是产出 $f(k^*)$ 与折旧 $\delta k^*$ 之间的缺口；它在 $k_g^*$ 处达到最大，那里资本的边际产出等于折旧率，$f'(k_g^*) = \delta$。}
\label{fig:macro-3}
\end{figure}

\begin{remark}[“等价”折旧]
黄金律条件 $f'(k_g^*) = \delta$ 中的 $\delta$ 应当被读作劳均资本被耗损的\textbf{等价}速率。在基准模型里它就是物理折旧 $\delta$。一旦加入人口增长，它就变成 $\delta + n$；再加入劳动增进型技术进步，它就变成 $\delta + n + g$。同一个条件——资本的边际产出等于等价折旧——自始至终成立；改变的只是我们往“等价折旧”里代入什么。
\end{remark}

\subsection{柯布-道格拉斯情形下的黄金律}

当生产为柯布-道格拉斯型时，黄金律储蓄率会呈现出一个格外干净的形式。令
\[
Y = A K^{\theta} L^{1-\theta}, \qquad 0 < \theta < 1,
\]
于是集约型生产函数为 $y = f(k) = A k^{\theta}$。

\ex[柯布-道格拉斯下的黄金律储蓄率]{
证明在 $Y = A K^{\theta} L^{1-\theta}$ 下，黄金律储蓄率为 $s_g = \theta$。
}
\sol{
集约型函数为 $f(k) = A k^{\theta}$，其边际产出为 $f'(k) = \theta A k^{\theta - 1}$。黄金律条件 $f'(k_g^*) = \delta$ 给出
\[
\theta A (k_g^*)^{\theta-1} = \delta
\quad\Longrightarrow\quad
\delta = \theta\,\frac{A (k_g^*)^{\theta}}{k_g^*} = \theta\,\frac{f(k_g^*)}{k_g^*}.
\]
代入储蓄率公式，
\[
s_g = \frac{\delta k_g^*}{f(k_g^*)} = \frac{\theta\, f(k_g^*)/k_g^* \cdot k_g^*}{f(k_g^*)} = \theta.
\]
}

$s_g = \theta$ 这个结果好记：黄金律储蓄率等于资本在产出中的份额。在竞争性经济里资本获得其边际产出的报酬，所以 $\theta$ 恰好是国民收入中归于资本的那个比例。这条规则说的是，把资本所赚取的那份收入精确地储蓄并再投资出去，不多也不少。储蓄超过资本份额就会过度积累——多出来的机器赚不回自身的损耗；储蓄少于资本份额则会让经济体达不到使消费最大化的存量。在这个意义上，黄金律要求在资本上的投资相对于资本对产出的贡献而言恰好划算（cost-effective）。

两项比较静态分析把这幅图补完整。更高的折旧率 $\delta$ 会拉低黄金律消费：经济体必须留出更多产出仅仅用来维持其资本，能享用的就更少了。而更高效的技术——每个 $k$ 上都有更高的 $f(k)$——会抬高黄金律消费，因为经济体生产得更高效；不过，与之相伴的黄金律储蓄率究竟是上升还是下降则是含糊的，要视 $f$ 的形状而定。

\section{储蓄率变化后的动态调整}

稳态告诉我们经济体最终落在哪里；它并没有告诉我们经济体如何走到那里。假设一位仁慈的社会规划者判断当前储蓄率过高——经济体已越过黄金律、过度积累了——于是把它调低到 $s_g$。消费、投资、产出与资本会随时间作何反应？答案完全取决于一个区分：哪些变量是\textbf{流量}、哪些是\textbf{存量}。

\state{存量缓慢调整，流量瞬间跳跃}{
\textbf{流量}（按单位时间度量的过程量，如消费 $c$、投资 $i$ 或产出 $y$）能在参数一改变的那一刻不连续地跳变。而\textbf{存量}（随时间积累起来的量，如资本 $k$）不能跳变；它只能通过投资这股流量被逐渐建立或耗减。追踪任何宏观调整，归根结底就是把这一区分理清楚，并搞清楚每个参数是怎样起作用的。
}

在储蓄率降到 $s_g$ 的那一刻，劳均资本 $k$ 还停在它原来那个更高的取值上——存量不能瞬间改变。由于 $k$ 暂时被钉住，产出 $y = f(k)$ 在冲击当下也不变。但那份固定产出的分配方式立刻就变了：投资 $i = s_g y$ 因 $s_g$ 变小而立即\textbf{下降}，消费 $c = (1 - s_g) y$ 则相应地立即\textbf{跳升}。两者都是流量，所以都在冲击当下移动；又因为产出在那一刻不变，投资的下降量恰好等于消费的上升量。

跳跃过后，缓慢的调整开始了。在更低的储蓄率下，投资如今不足以补足折旧，$s_g f(k) < \delta k$，故 $\Delta k < 0$，劳均资本朝新的、更低的稳态漂落。随着 $k$ 下降，产出 $y = f(k)$ 也随之下降，并紧跟着 $k$，因为 $y$ 通过生产函数与 $k$ 绑定在一起。随着产出下降，消费与投资——如今都是一份不断缩小的产出的固定比例——也逐渐随之下滑。消费在冲击当下纵身跃起之后，又缓缓回落到它新的稳态水平，但这一水平仍然高于此前（这正是迁向黄金律的初衷）。图~\ref{fig:macro-4} 展示了这四条时间路径：流量 $i$ 与 $c$ 在政策变动当期的陡然不连续跳跃、存量 $k$ 的平滑下滑，以及紧跟着 $k$ 的产出 $y$。

\begin{figure}[h]
\centering
\begin{tikzpicture}[scale=1.0]
  % ===== shared panel geometry =====
  % each panel is 5 wide, 3.6 tall; the jump happens at t = tj
  \def\W{5.2}\def\H{3.6}\def\tj{0.9}
  % horizontal offsets of the two columns and vertical offsets of the two rows
  \def\xR{7.0}   % x-shift of the right column
  \def\yB{-5.0}  % y-shift of the bottom row

  % ---------- helper macro: one panel ----------
  % #1 = (dx,dy) shift   #2 = vertical-axis variable label
  % drawn relative to (0,0); axes go up to (\W,0) and (0,\H)

  % ============ TOP-LEFT : k ============
  \begin{scope}[shift={(0,0)}]
    \draw[-Latex] (0,0) -- (\W,0) node[below right] {$t$};
    \draw[-Latex] (0,0) -- (0,\H) node[left] {$k$};
    \def\kstar{2.85}\def\kg{1.15}
    % asymptote
    \draw[densely dashed] (0,\kg) -- (\W,\kg);
    \node[left] at (0,\kg) {$k_g^{*}$};
    \node[left] at (0,\kstar) {$k^{*}$};
    % flat at k*, then smooth decline toward k_g*
    \draw[line width=1.1pt] (0,\kstar) -- (\tj,\kstar);
    \draw[line width=1.1pt,smooth,samples=120,domain=\tj:\W]
      plot (\x,{\kg + (\kstar-\kg)*exp(-1.05*(\x-\tj))});
  \end{scope}

  % ============ TOP-RIGHT : y ============
  \begin{scope}[shift={(\xR,0)}]
    \draw[-Latex] (0,0) -- (\W,0) node[below right] {$t$};
    \draw[-Latex] (0,0) -- (0,\H) node[left] {$y$};
    \def\ystar{2.85}\def\yg{1.15}
    \draw[densely dashed] (0,\yg) -- (\W,\yg);
    \node[left] at (0,\yg) {$y_g^{*}$};
    \node[left] at (0,\ystar) {$y^{*}$};
    \draw[line width=1.1pt] (0,\ystar) -- (\tj,\ystar);
    \draw[line width=1.1pt,smooth,samples=120,domain=\tj:\W]
      plot (\x,{\yg + (\ystar-\yg)*exp(-1.05*(\x-\tj))});
  \end{scope}

  % ============ BOTTOM-LEFT : c ============
  \begin{scope}[shift={(0,\yB)}]
    \draw[-Latex] (0,0) -- (\W,0) node[below right] {$t$};
    \draw[-Latex] (0,0) -- (0,\H) node[left] {$c$};
    \def\cstar{1.30}\def\cg{1.70}\def\cpk{2.95}
    % asymptote c_g* sits ABOVE the starting level c*
    \draw[densely dashed] (0,\cg) -- (\W,\cg);
    \node[left] at (0,\cg) {$c_g^{*}$};
    \node[left] at (0,\cstar) {$c^{*}$};
    % flat at c*, then JUMP UP to the peak, then decline toward c_g*
    \draw[line width=1.1pt] (0,\cstar) -- (\tj,\cstar);
    \draw[line width=1.1pt] (\tj,\cstar) -- (\tj,\cpk);
    \draw[line width=1.1pt,smooth,samples=120,domain=\tj:\W]
      plot (\x,{\cg + (\cpk-\cg)*exp(-0.78*(\x-\tj))});
  \end{scope}

  % ============ BOTTOM-RIGHT : i ============
  \begin{scope}[shift={(\xR,\yB)}]
    \draw[-Latex] (0,0) -- (\W,0) node[below right] {$t$};
    \draw[-Latex] (0,0) -- (0,\H) node[left] {$i$};
    \def\istar{2.85}\def\ig{0.95}\def\idr{2.05}
    % asymptote i_g* below
    \draw[densely dashed] (0,\ig) -- (\W,\ig);
    \node[left] at (0,\ig) {$i_g^{*}$};
    \node[left] at (0,\istar) {$i^{*}$};
    % flat at i*, then JUMP DOWN to a level, then decline toward i_g*
    \draw[line width=1.1pt] (0,\istar) -- (\tj,\istar);
    \draw[line width=1.1pt] (\tj,\istar) -- (\tj,\idr);
    \draw[line width=1.1pt,smooth,samples=120,domain=\tj:\W]
      plot (\x,{\ig + (\idr-\ig)*exp(-0.95*(\x-\tj))});
  \end{scope}
\end{tikzpicture}

\caption{储蓄率被砍到 $s_g$ 之后，资本 $k$、产出 $y$、消费 $c$ 与投资 $i$ 的时间路径。流量 $i$ 与 $c$ 在冲击当下不连续地跳跃（投资下降、消费上升）；随后存量 $k$ 缓慢下降，拖着 $y$、$c$、$i$ 一道朝新的稳态走低。}
\label{fig:macro-4}
\end{figure}

这条路径还有一个值得点名的特征：向新稳态的收敛是缓慢而减速的。随着 $k$ 接近它新的歇脚点，投资与折旧之间的缺口缩小，故 $\Delta k$ 缩小，经济体越来越轻地朝稳态蹭去。抵达新的均衡是许多期才能完成的工作，而非一蹴而就。

\section{扩展的索洛模型}

基准模型让劳动力恒定、并假设没有技术变化。这两点都明显与事实不符，而放松它们正是让模型能够谈论持续增长的关键。指导原则就是动态框架中所强调的那条：新的因素只通过\textbf{运动方程}进入；静态关系并不改变。在伸手去拿代数工具之前，先把直觉想清楚。

\subsection{人口增长}

设劳动力以恒定速率 $n$ 增长，
\[
\frac{\Delta L}{L} = n.
\]
总量的运动方程毫发未动，$K' = (1-\delta) K + I$。改变的是向劳均形式的换算，因为工人数本身如今在上升。记 $k = K/L$，并除以下一期更大的劳动力，
\[
k'(1 + n) = (1-\delta) k + i.
\]
减去 $k$ 并化简，
\[
\ba
\Delta k = k' - k
&= i - \delta k - n k' \\
&= i - \delta k - n k + n\!\pr{\frac{K}{L} - \frac{K'}{L'}} \\
&\approx i - (\delta + n) k,
\ea
\]
其中 $n\,(K/L - K'/L')$ 一项是两个微小变化之积，在一阶近似下可以忽略。于是我们得到人口增长情形下的核心结果，
\begin{equation}
\Delta k \approx s f(k) - (\delta + n) k.
\label{eq:solow-popn}
\end{equation}

其解读立竿见影，而且恰是直觉所预言的那样。与基准模型相比，劳均资本被耗损的等价速率从 $\delta$ 升到了 $\delta + n$。除了替换掉用坏的机器（$\delta k$ 一项），经济体如今还必须\textbf{为每期加入劳动力的新工人配备资本}（$n k$ 一项），仅仅为了不让劳均资本下降。说到底，人口增长起的作用就像额外的折旧。稳态条件变为 $s f(k^*) = (\delta + n) k^*$，黄金律条件变为 $f'(k_g^*) = \delta + n$。

\begin{remark}[人口增长、马尔萨斯陷阱与 CRS 的局限]
由于资本的回报受边际报酬递减支配，而人口增长又不断稀释着每个工人所拥有的资本，更快的人口增长会拉低稳态劳均资本、从而拉低劳均产出——人人都更穷了。推到极致，这就是\textbf{马尔萨斯陷阱}（第~\ref{ch:malthus}~章）的逻辑：收入的任何增益都被更庞大的人口耗散掉。这一机制稳稳地立在规模报酬不变的假设之上。现实中人的群集可以产生正的外部性——思想、专业分工、稠密的市场——以至于社会可能呈现\textbf{规模报酬递增}（IRS）。况且，生育是一种主动的选择，把人口增长当作一个内生的决策、而非一个外生的参数来处理，才是更令人满意的建模立场，而这一立场正是索洛框架所不采取的。
\end{remark}

\subsection{劳动增进型技术进步}

要产生生活水平的持续增长，我们必须加入技术进步。索洛是通过一个刻意特殊化的装置来做这件事的：技术作为劳动的乘数进入生产函数，
\[
Y = F(K, E \cdot L),
\]
其中 $E$ 度量\textbf{劳动效率}（efficiency of labor），并以恒定速率 $g$ 增长。乘积 $E \cdot L$ 就是\textbf{有效工人}（effective workers）的供给。我们通常把技术进步想象成不断上升的全要素生产率、改善着所有投入的使用；而把它建模为\textbf{劳动增进型}（labor-augmenting）则是一个便于处理的选择，它让经济体能够安顿到一条干净的平衡路径上。（如果技术也增进资本，那就相当于直接改进全要素生产率了。）

运动方程依旧不变，$K' = (1-\delta)K + I$。现在我们把一切按\textbf{有效}工人来度量，定义 $k = K/(EL)$、$i = I/(EL)$。由于有效劳动以 $n + g$ 的速率增长（人口以 $n$、效率以 $g$），把运动方程除以下一期的有效劳动，
\[
k'(1+n)(1+g) = (1-\delta) k + i.
\]
展开并丢掉二阶项 $ng$，
\[
k'(1 + n + g) = (1-\delta) k + i,
\]
经由与之前相同的步骤整理后得到
\begin{equation}
\Delta k \approx s f(k) - (\delta + n + g) k.
\label{eq:solow-tech}
\end{equation}
如今等价折旧率是 $\delta + n + g$：每个有效工人的资本被物理折旧耗损、被人口增长稀释、又被劳动效率的增长再稀释一次。在稳态处令 $\Delta k = 0$ 得到
\[
(\delta + n + g)\, k^* = s f(k^*) = i^*.
\]

关键的问题在于，这对真实的生活水平意味着什么——而生活水平是按\textbf{人均}、而非按有效工人来度量的。在稳态处 $k^* = K/(EL)$ 不变，所以\textbf{人均}资本与产出为
\[
\frac{K}{L} = k^* E = k^* E_0 (1+g)^t,
\qquad
\frac{Y}{L} = y^* E = y^* E_0 (1+g)^t.
\]
尽管把 $g$ 加进等价折旧会拉低以有效工人为单位度量的稳态 $k^*$（从而拉低 $y^*$），但乘性因子 $E_0(1+g)^t$ 无界增长，完全占据主导。于是人均产出与人均资本永远以 $g$ 的速率增长。这是这个模型的核心成就：\textbf{劳动增进型技术进步是人均收入持久增长的唯一来源}。与此同时，总产出 $Y$ 等于 $EL$ 乘以它的有效工人取值，故以 $n + g$ 的速率增长。

\begin{remark}[作为公共品的技术]
为什么劳动增进型技术抬高的是所有人的收入，而不是自己索取一份要素报酬？模型中的关键在于，效率 $E$ 是\textbf{非竞争性的}（non-rivalrous）：我用一个想法，并不妨碍你也用它。资本与劳动是竞争性要素，在要素市场上交易、赚取回报；而技术在这个意义上并不是一种要素，而是一种公共品，其好处归于全体劳动者。我们可以把 $E$ 想象成一个乘数，让社会用 $L$ 个人就实现本来需要 $E \cdot L$ 个人才能达到的成效；当我们盘点人类福祉时，又把这 $E$ 个有效工人合并回 $L$ 个真实的人，于是增益就以人均收入的不断上升显现出来。现实中专利、保密或其他摩擦使技术部分地变得可排他，结果便会与这个无摩擦的模型有所偏离。
\end{remark}

\subsection{稳态增长率：小结}

模型的三个版本之间的差别，仅在于长期究竟什么在增长。集约型变量（按有效工人）在每个稳态里都按构造保持不变；差别在于人均与总量的增长率，汇总于表~\ref{tab:solow-summary}。

\begin{table}[h]
\centering
\caption{索洛模型三个版本在稳态下各变量的长期增长率。}
\label{tab:solow-summary}
\begin{tabular}{@{}lccc@{}}
\toprule
稳态下 & 基准 & 人口增长 & 人口与技术增长 \\
\midrule
按有效工人 & 不变 & 不变 & 不变 \\
人均 & 不变 & 不变 & $g$ \\
总量 & 不变 & $n$ & $n + g$ \\
\bottomrule
\end{tabular}
\end{table}

这张表把上文的讨论一一读出。当劳动与技术都不增长时，一切都是平的。当只有人口增长时，长期来看人均产出依旧是平的——经济体更大了，但人均上并没有更富——而总产出以 $n$ 增长。只有当劳动增进型技术进步在场时，人均收入才会增长，速率为 $g$，总量则以 $n + g$ 增长。

\section{关于索洛模型的更多讨论}

索洛框架尽管简单，却支撑得起一组丰富的进一步观察。我们把最有用的几条收集在这里。

\subsection{经济增长的三个阶段}

从经验上看，各国似乎要经历三个大的阶段，每个阶段由一种不同的增长机制支配。

\begin{enumerate}
    \item \textbf{贫困陷阱}（人均 GDP 约低于 \$3{,}000）。劳均资本极低，所以资本的边际产出——从而投资的回报——极高。处于这一阶段的经济体可以仅靠积累资本就快速增长，这个过程称为\textbf{资本深化}（capital deepening）。诸如中国这样的快速追赶式增长，恰恰是基准索洛模型对一个资本稀缺经济体所作的预言，而非什么反常现象。
    \item \textbf{中等收入转型阶段。}随着资本深化，报酬递减开始显现，靠积累带来的轻松增长逐渐消退。要继续增长，经济体必须从要素积累转向创新驱动的增长。这一转型很难——转型失败就是所谓的\textbf{中等收入陷阱}——驾驭它通常需要持续的制度改革。
    \item \textbf{高收入阶段。}这里增长压倒性地来自技术创新，因为进一步的资本深化已收效甚微。在这个区制里，是 $g$ 这一项、而非储蓄率，在驱动生活水平。
\end{enumerate}

\subsection{增长核算与索洛余值}

为了度量增长究竟来自何处，设产出沿着一条带全要素生产率 $z$ 的柯布-道格拉斯路径走，
\[
Y = z\, K^{\alpha} L^{1-\alpha}.
\]
取对数再微分，得到\textbf{增长核算}（growth accounting）分解，
\begin{equation}
\frac{\d Y}{Y} = \frac{\d z}{z} + \alpha\,\frac{\d K}{K} + (1-\alpha)\,\frac{\d L}{L}.
\label{eq:growth-acct}
\end{equation}
产出增长是生产率增长、加上资本增长与劳动增长各自贡献之和，每项都以其收入份额为权重。历史上 $\alpha \approx 1/3$。化为劳均形式，$y = z k^{\alpha}$，分解坍缩为
\[
\frac{\d y}{y} = \frac{\d z}{z} + \alpha\,\frac{\d k}{k}.
\]

$z$ 一项就是\textbf{索洛余值}（Solow residual）：把资本与劳动可度量的贡献扣除之后，产出增长中剩下的那部分。它之所以称为余值，正是因为它无法被直接观测，而是从式~\eqref{eq:growth-acct} 中反推出来的，$\d z / z = \d Y/Y - \alpha\, \d K/K - (1-\alpha)\, \d L/L$。

\defn{索洛余值}{
\textbf{索洛余值} $z$（常被称为全要素生产率）是产出中不能由资本与劳动的可度量数量所解释的那部分。它捕捉了在给定投入下抬高产出的一切，包括
\begin{itemize}
    \item 人力资本（凝结在工人身上的教育与技能）；
    \item 技术进步；
    \item 外部性（公共品的溢出，与那种可排他、收费的物品相反）；
    \item 制度质量（厂商的组织与管理知识、专利保护、对腐败的治理）。
\end{itemize}
}

请注意它与上一节劳动增进型 $E$ 的对照：$E$ 只改善劳动，而余值 $z$ 乘的是整个生产函数，因此同时抬高资本与劳动的生产率。理解 $z$ 的一个有用方式是，它本身是一种\textbf{结果}——是对教育、研究与制度的投资所形成的均衡产物——所以人们可以用分析任何其他均衡对象时所用的同一套供求逻辑，去分析它的水平与增长。促进增长的政策杠杆，直接从余值的内容里读出来：提高储蓄率、投资人力资本、鼓励技术进步、建立正确的制度。

\subsection{平衡增长路径与收敛}

扩展模型的稳态是一条\textbf{平衡增长路径}（balanced growth path，BGP）：沿着这条轨迹，产出、资本与消费全都以恒定（且彼此一致）的速率增长。和基准稳态一样，BGP 是一个吸引子——若经济体被撞离它，各变量会随时间收敛回来。

两个生产率度量有助于在路径上给经济体定位。\textbf{劳动生产率}与\textbf{资本生产率}分别为
\[
\frac{Y}{L} = z\!\pr{\frac{K}{L}}^{\!\alpha},
\qquad
\frac{Y}{K} = z\!\pr{\frac{K}{L}}^{\!\alpha - 1}.
\]

\begin{remark}[劳动生产率不是工资]
劳动生产率 $Y/L$ 不应与平均工资相混淆，因为这里的 $L$ 是\textbf{在业}工人数，而非整个人口；失业、劳动参与率与人口结构在劳均产出与人均收入之间打入了一道楔子。
\end{remark}

索洛逻辑现在交出它最锋利的一课。资本-劳动比 $K/L$ 的增长\textbf{不}可持续——报酬递减保证了它会逐渐熄火；而余值 $z$ 的增长\textbf{是}可持续的，且它同时抬高劳动生产率与资本生产率。成熟经济体主要靠不断上升的全要素生产率增长；比如今天的美国，其增长大半要归功于余值、而非资本深化。

最后，模型对收敛有一个重要的限定。基准模型那种无条件收敛——人人最终一样富——只在各经济体共享相同的基本面时才成立。由于各国的制度与技术路径不同，经济体收敛到的并不是单独一个共同的稳态，而是收敛到\textbf{各自}的平衡增长路径。那些共享相似制度与技术的国家，确实会收敛到一条共同的路径，这一现象被称为\textbf{条件收敛}或\textbf{俱乐部收敛}（conditional / club convergence）。技术进步的内生性——$g$ 本身取决于政策、制度与研究这一事实——恰恰是索洛模型留而未解的，而那正是后来的增长理论的起点。

\state{索洛能解释什么、不能解释什么}{
索洛模型解释了为什么资本贫乏的经济体先快速增长而后放慢、为什么资本积累无法维持增长，以及为什么生活水平的长期增长必须来自技术进步。它不能解释那技术进步从何而来，也不能解释为什么有些国家的技术与制度改善得比另一些更快。它并不能解释一切——但它能解释相当多的东西，并且通过用优化的家庭取代索洛那条机械的储蓄规则，它框定了新古典模型（第~\ref{ch:neoclassical}~章）将要接手的问题。
}
